%% The following is a directive for TeXShop to indicate the main file
%%!TEX root = diss.tex
\chapter{Discussion on Random Walk Model Generalizations}

%================================
\section{Introduction}
%================================

In this chapter we'll push the operator notation further to see what kinds of random walk models have a relatively easy functional equation to write. We start with the following definition for a random walk model in $\ze^d$.
\begin{definition} Let $\mc{D}_0 \subseteq \mc{D}$ be regions in $\ze^d$ and $S$ be a set of $d$-tuples in $\ze^d$. A triple $\mc{M} = (\mc{S}, \mc{D},\mc{D}_0)$ is a random walk model $\mc{M}$ with steps $\mc{S}$, living in the region $\mc{D}$ and interacting with region $\mc{D}_0$. We may impose a starting location later on.
\end{definition}
In other words, we are given a set of steps $\mc{S}$ and a starting location in $\mc{D}$. Then we want to evolve the random walks restricted to the region $\mc{D}$. We also have some special locations $\mc{D}_0$ within $\mc{D}$ and we wish to count the interactions of the random walks with points in this region. The interactions may be further specified later on.
\begin{definition} Let $\mc{M}$ be a random walk model and $Q$ its associated generating function counting length with $t$, ending locations with $\row xd$ and interactions with $\mc{D}_0$ with some multi-index variable $a$. We write
\begin{align}
	Q = Q(\row xd,a;t) = \sum_{k=0}^{\infty} q_k(\row xd, a)t^k
\end{align}
where $q_k$ are Laurent poynomials in the $x$ and $a$ variables. 
\end{definition}
\begin{definition} Suppose $D$ is an operator involving the $x$ and $a$ variables but not $t$. We shall apply $D$ on power series $F = \sum_{k=0}^\infty f_k t^k$ in the following linearized way
\begin{align}
	DF = \sum_{k=0}^\infty D(f_k) t^k
\end{align}
to avoid issues of expanding $F$ in Laurent series in some of the $x$ variables.
\end{definition}

\begin{definition} Let $\mc{M}$ be a given model. We'll denote with $S$ the Laurent polynomial (or series) keeping inventory of the stes $\mc{S}$. Let $K = 1-tS$ be the Kernel of $\mc{M}$. Let $E$ be an operator that eliminates monomials $x_{[d]}^P$ when $P$ is not in $\mc{D}$. Let $L$ be an operator that marks (or labels) monomials $x_{[d]}^P$ for $P \in \mc{D}_0$ in a specified way in the model. Note that $L$ is invertible with inverse $L^{-1}$ that is given by marking with $1/a$ instead.
\end{definition}

This operator notation allows us to write down functional equations for $Q$ really quickly. We have
\begin{theorem} Let $\mc{M}$ be a random walk model and $Q$ its associated generating function with an initial term $q_0 = q$. Then $Q$ satisfies the following functional equations
\begin{align}
	(1-tLES)Q &= q\\
	(L^{-1}-tES)Q &= L^{-1}(q) = \alpha
\end{align}
\end{theorem}
\begin{proof} The second functional equation is obtained by applying $L^{-1}$ on the first one. It only remains to show that the first functional equation holds. We do this via iteration. We start with $Q$ and multiply it with $tS$. This gives $tSQ$ or an enumeration of all walks extended by a single step. Here the power of $t$ is raised by one to keep track of this extension. Then we apply $E$ to this result to eliminate all random walks that venture outside of the region $\mc{D}$. Finally we mark those walks that have arrived at locations in $\mc{D}_0$ by applying $L$ to $E(tSQ)$. The resulting expression is almost $Q$ but with one term missing, namely $q_0 = q$. That is we have transformed $Q$ to itself in a nontrivial way giving us the functional equation
\begin{align}
	Q = q + LE(tSQ) = q + tLE(SQ)
\end{align}
which we re-write in the form at the start of the theorem.
\end{proof}
\begin{corollary} We have the following kernel forms
\begin{align}
	KQ &= q - t(1-LE)SQ\\
	KQ &=\alpha + (1-L^{-1})Q -t(1-E)SQ.
\end{align}
If there is no extraction involved we have the kernel forms
\begin{align}
	KQ &= q - t(1-L)SQ\\
	KQ &=\alpha + (1-L^{-1})Q.
\end{align}
If there is no marking involved then we have the simple kernel form
\begin{align}
	E(KQ) = q.
\end{align}
\end{corollary}
 The kernel forms are obtained as before since it was mostly a symbollic calculation. We use the fact that $EQ = Q$ since $Q$ is already composed of terms living in $\mc{D}$. We won't consider more general functional equation statements in this work. It is possible to also consider step sets $\mc{S}$ varrying with location as we saw in the honey-comb lattice chapter but writing a functional equation for this type of problem will be left as a (rather long) exercise to the reader.

The second form of the functional equation involving $L^{-1}$ should have less terms than the first one by virtue of separating $L$ and $E$ from one another. Using extraction operators we can constrain random walks to lie within a finite region, semi-infinite strip, a semi-infinite sector or any other discrete shape (with diminishing returns for more complicated shapes of course). Furthermore using marking operators we can easily track the behavior of random walks as they cross or interact with boundaries or interact with regions within. For example we may consider walks interacting with the boundary of a box within some region in space. 

\begin{example} We can also consider a different approach where instead of removing walks that venture outside of $\mc{D}$ we just mark them with some variable $z$ for later deletion. Suppose we are given some $\mc{M}$ and we want to frame in this new way. We have a new random walk model $\wcal{M}$ with $\wcal{D} = \ze^d$ and $\wcal{D}_0 = \mc{D}_0 \cup \mc{D}^c$ where we mark monomials in $\mc{D}^c$ with $z$ and those in $\mc{D}_0$ as in $\mc{M}$. Here $\widetilde{E}$ is the identity since we don't remove any walks. We have the functional equations
\begin{align}
(1-t\widetilde{L} S)\widetilde Q &= q\\
(\widetilde{L}^{-1}-tS) Q &= \widetilde{\alpha}
\end{align}
and kernel forms
\begin{align}
	K\widetilde Q &= q - t(1-\widetilde L)S \widetilde Q\\
	K\widetilde Q &=\widetilde \alpha + (1-\widetilde{L}^{-1}) \widetilde Q.
\end{align}
where $\widetilde Q = Q + O(z)$ wher $Q$ is the generating function associated with model $\mc{M}$. We just need to extract $[z^0]$ from $\widetilde Q$ to obtain $Q$. 
\end{example}

We can talk about some of the abstract properties of the operators $L$ and $E$. The operator $E$ has the eigenvalues 1 or 0 for monomials $x_{[d]}^P$ where $P \in \mc{D}$ or $P \in \mc{D}^c$ respectively. The latter characterizes the monomials contained in $\ker E$. Similarly $L$ has eigenvalues 1 for when $P \in \mc{D}_0^c$ and the marking monomial $a$ for when $P\in \mc{D}_0$. Here $a$ will varry from point to point. Also $L^{-1}$ will have the reciprocal eigenvalues of $L$.

The task for the rest of the chapter is to come up with representations for the operators involved. In general this should not be an easy task because if we could express arbitrary operators $L$ and $E$ with a small number of coefficient extractions we'd also be able to use this scheme to compress data without loss, which is impossible in general.

%================================
\section{Marking Operators}
%================================
The example in the previous section shows that we can talk about marking operators on their own without considering operators $E$. Alternatively, the operators $E$ are just marking operators which apply a label $a = 0$. Therefore we can combine the discussions for operators $L$ and $E$ into one holistic approach. Whenever invertibility is concerned the discussion will apply to operators $L$ only. Suppose that we wish to mark points $p \in P\subseteq \ze^d$ with labels $a_p$ via an operator $\Gamma_P$. It can be written as
\begin{align}
	\Gamma_P = 1 + \sum_{p\in P}(a_p-1) \ext{x_{[d]}^p}
\end{align}
The form $a_p-1$ follows from the fact that we have the original term and we need to remove it before marking with $a_p$. In the case where we mark points for deletion $P = {\mc{D}}^c$ and $a_p = a$. The good thing is that in a lot of cases we don't need all the terms in the operator and in some cases we even find the operator has a portion that factors. We just need to ensure that we mark all the walks that touch the boundaries or go within the region marked for deletion. For example we can take the walks in one dimensions on the non-negative axis with steps of size at most $n$ in the negative direction. Then we can demonstrate the truncation with
\begin{align}
	\Gamma = 1 + (a-1)\sum_{k=1}^{\infty} \ext{\ol{x}^k} = 1 + (a-1)\sum_{k=1}^{n} \ext{\ol{x}^k}
\end{align}
We only need to consider how far a walk can step outside of the region $\mc{D}$ via the step set in $S$. In two dimensions we have
\begin{align}
	\Gamma = 1 + (a-1)\sum_{m\text{ or }n < 0}^{\infty} \ext{\ol{x}_1^m \ol{x}_2^n}
\end{align}
Without $a$ this reduces to
\begin{align}
	\Gamma = 1- \ext{\ol{x}_1}- \ext{\ol{x}_2}+ \ext{\ol{x}_1\ol{x}_2} = (1-\ext{x}_1)(1-\ext{x}_2)
\end{align}
in which case we also see that the operator can factor into two or more operators sometimes. When the boundary is complicated we may not find an a finite form for $\Gamma$ let alone see any factoring. However, there are some nice special cases where we can say a lot about the operator $\Gamma$. We can build the theory. We have only finitely many steps (and sizes of steps) in $S$. We say $\Gamma$ is finite if it can be represented by finitely many coefficient extraction terms. Suppose that $\mc{D}$ is a finite box in $\ze^d$. Then $\Gamma$ is finite. To see this suppose that $B$ has $x_k$ coordinates in the interval $[a_k,b_k]$ and that stepping outside from $B$ from within $B$ via step polynomial $S$ leads to $x_k$ coordinates within the interval $[m_k,n_k]$. Then $\Gamma$ can be written finitely with extraction part
\begin{align}
	\Gamma_{B^c} = \sum_{k=1}^{d} \left(\sum_{i= m_k}^{a_k-1}\ext{x_k^i}+\sum_{i= b_k+1}^{n_k}\ext{x_k^i}\right)
\end{align}
 The discussion continues for semi-infinite and infinite boxes with finite marked regions within. The analysis is the same as above except that now one or more of $a_k, b_k$ may be infinite, in which case we actually lose parts of $\Gamma_{B^c}$ and the operator is still finite. We should see if we can do regions whose complements are boxes and then build up to finite regions and intersections of finite regions with boxes.
 
\begin{definition} A coefficient extraction operator $\Gamma$ in the variables $x_1,\ldots,x_d$ is of finite type if it can be written in a form with only finitely man terms like
\begin{align}
	\Gamma = \sum_{\substack{I\subseteq [d]\\ J \subseteq \ze^d}}a_{I,J} x_I^J[x_I^J]
\end{align}
where each $|J| < \infty$. Then $\Gamma$ is of infinite type if it's not of finite type. We'll just say $\Gamma$ is finite or infinite.
\end{definition}

\begin{definition} A random walk model with finitely many steps is of finite type if its extraction and marking operator $\Gamma$ is of finite type and its starting condition $q$ is of finite degree. If not, the model is of infinite type.
\end{definition}
Not surprisingly all the models so far has been of finite type. In the next section we'll see the model of self-avoiding walks is most likely of infinite type.

\begin{theorem} The generating function $Q$ for a random walk model of finite type is D-finite in the variables $x_1,\ldots,x_d$.
\end{theorem}
\begin{proof} We'll use the kernel form for the functional equation
\begin{align}
	KQ + t(1- \Gamma)SQ = q.
\end{align}
Here $\Gamma' = 1- \Gamma$ is purely a coefficient extractor and is of finite type. We'll prove this by finding the differential equation explicitly. Since the model is of finite type we can bound the degree of $x_k$ that we see on either side of the functional equation. In one scheme we let $A_k$ be the largest power of $x_k$ that we need to clear a negative power of $x_k$ in $\Gamma'$ and in $q$. Then we let $B_k$ be the largest power of $x_k$ that we see on either side. Once we multiply by $x_k^{A_k}$ we see that the degree of $x_k$ in the equation is $A_k+B_k$. We take the functional equation and multiply both sides by $x_{[d]}^{A}$
\begin{align}
	x_{[d]}^{A}(KQ + t\Gamma'SQ) = x_{[d]}^{A}q
\end{align}
and we obtain a polynomial equation in the variables $x_1,\ldots,x_d$. We take the derivative operator
\begin{align}
	D = \del_{[d]}^{A+B+\one} = \del_1^{A_1+B_1+1}\ldots\del_d^{A_d+B_d+1}
\end{align}
and apply it to both sides of the equation to get
\begin{align}
	D(x_{[d]}^{A}(KQ + t\Gamma'SQ)  = D(x_{[d]}KQ) = D(x_{[d]}^{A}q) = 0
\end{align}
which follows since we're taking more derivatives than the degrees of the monomials which clears all of them. The other scheme lets us obtain a potentially lower degree differential equation at the cost of it beind inhomogeneous. Let $A_k$ be the power of $x_k$ that we need to clear denominators on the LHS of the functional equation and $B_k$ be the largest power of $x_k$ on that side. Then we have $D$ as before but this time it may have lower $A$ and $B$ values since we are not taking $q$ into account. We have
\begin{align}
	D(x_{[d]}^{A}(KQ + t\Gamma'SQ)  = D(x_{[d]}KQ) = D(x_{[d]}^{A}q)
\end{align}
\end{proof}
For example let the powers of $x_k$ be $\{-3,-1,4\}$ and $\{-2,0,6\}$ on the LHS and RHS respectively. We see that $A_k = 3$ and $B_k = 6$. Then we're taking the $A_k + B_k + 1 = 10$-th derivative of both sides of the functional equation. In The second scheme we take $A_k = 3$ still but $B_k = 4$ and we are taking the 8-th derivative w.r.t. $x_k$ this time. Even in the quarter-plane we have models that have a generating function that is not D-finite (with respect to the $t$ variable) but it turns out that every model so far considered is D-finite in the $x_1,\ldots,x_d$ variables.

In order to solve the differential equation we need to investigate the kernel of the differential operator.
\begin{corollary}
The generating function $Q$ of a finite random model is given by
\begin{align}
	x_{[d]}^{A}KQ(x_1,\ldots,x_d;t) = \sum_{i=1}^d \sum_{j=0}^{B_i}Aij(x_1,\ldots,x_{i-1},x_i,\ldots, x_d,t) x_i^j
\end{align}
where $A_{ij}$ are functions that need to be determined.
\end{corollary}
We could have also obtained this kind of solution from going through the calculations in the functional equation and extracted coefficients. However, this method is more efficient. If we had used the second scheme in the proof we would have ended up with a slightly different form but maybe it would have had a slightly lower degree. It would involve an indefinite integral of $q$ and elements in the kernel of the differential operator.

%================================
\section{Self-avoiding walks}
%================================

We can handle self-avoiding random walks with this theory. If we pick a marking operator $\Gamma$ that marks every location on the plane this will allow us to track the random walks' history. All we'll need to do in the end is remove those walks which have stepped at some location more than once. However we'll do this problem in a bit more detail as it serves as an example of the theory. To use the generating function iteration properly we need a sort of limiting argument. For the first $n$ steps we "only" require an operator covering the square region of $(2n+1)^2$ lattice points. The overall operator is defined through a limiting process of applying the full operator (or truncated version) on a finite Laurent polynomial. Let
\begin{align}
\Gamma_{i,j} = 1+b_{i,j}\ext{x_1^ix_2^j}
\end{align}
where as usual $b_{i,j} = a_{i,j}-1$. This operator takes a generating function and marks the monomial $x_1^ix_2^j$. We'll be using a lot of these together so we let
\begin{align}
\Gamma_n = \prod_{-n\leq i,j\leq n}\Gamma_{i,j}
\end{align}
To find a functional equation for SAWs on the plane we perform an iteration. Let $Q$ be the generating function for walks with history, that is they have a mark for each coordinate they pass through and let $Q^{\leq n}$ be the truncated version of this generating function. It's given by
\begin{align}
Q^{\leq n} = \sum_{k=0}^{n} Q_k t^k
\end{align}
where the initial condition is $Q_0 = a_{0,0}$. Starting from $Q_n$ we can find $Q_{n+1}$ by iterating with the step set $S$, marking with $\Gamma_{n+1}$ as follows
\begin{align}
	Q_{n+1} = \Gamma_{n+1}SQ_n
\end{align}
Using a computer we can easily iterate this to get more and more steps. We can also get $Q^{\leq n+1}$ from $Q^{\leq n}$ since
\begin{align}
	t\Gamma_{n+1}SQ^{\leq n} & = \sum_{k=0}^n \Gamma_{n+1}SQ_kt^{k+1}\\
	& = \sum_{k=0}^n \Gamma_{k+1}SQ_kt^{k+1}\\
	& = \sum_{k=0}^n Q_{k+1}t^{k+1}\\
	&= Q^{\leq n+1}-Q_0
\end{align}
or in other words
\begin{align}
	 Q^{\leq n+1} = a_{0,0} + t\Gamma_{n+1}SQ^{\leq n}
\end{align}
In the interest of finding a functional equation we would like to take the limiit of this as $n \to \infty$. For this we define $\Gamma$ to be the operator marking all monomials $x_1^ix_2^j$ and acting in the following way
\begin{align}
	\Gamma x^i y^j = \Gamma_{\max(|i|,|j|)} x^i y^j
\end{align}
and extended by linearity. The limiting process is
\begin{align}
	Q = a_{0,0} + t \Gamma SQ
\end{align}
We finally have the functional equation for random walks with full history
\begin{align}
	(1 -t\Gamma S)Q = a_{0,0}.
\end{align}
From this we can obtain the generating function $W$ for the SAWs by removing all walks with power of mark $a_{i,j}$ greater than 1. For this we need a removal operator. Let
\begin{align}
	R_{i,j} &= \ext{a_{i,j}^0} + \ext{a_{i,j}}\\
	R_n &= \prod_{-n\leq i,j \leq n} R_{i,j}\\
	R &= \lim_{n\to \infty} R_n
\end{align}
where the limit is understood as in the example of $\Gamma$ above. The solution for $Q$ can also be restated as
\begin{align}
	Q = \sum_{k=0}^{\infty} t^k (\Gamma S)^{k}(a_{0,0}) = \sum_{k=0}^{\infty} t^k (\Gamma_k S)^{k}(a_{0,0}) 
\end{align}
We can either redo the analysis applying $R$ at every iteration or we can apply $R$ to $Q$ and get
\begin{align}
	W = RQ = \sum_{k=0}^{\infty} t^k R(\Gamma S)^{k}(a_{0,0}) = \sum_{k=0}^{\infty} t^k (R\Gamma S)^{k}(a_{0,0}) 
\end{align}
We can write this as a functional equation
\begin{align}
	\mc{L} W = (1-t\Sigma)W = (1-tR\Gamma S)W = a_{0,0}
\end{align}
%================================
\subsection{Eigenpairs}
%================================
We can try finding eigenvalues and eigenfunctions of $\mc{L}$. Let $\mc{L}u = \lambda u$. Writing it out leaves us with
\begin{align}
	t\Sigma u = (1-\lambda)u
\end{align}
So if $\lambda = 1$, $u\in \ker(\Sigma)$. That is $R\Gamma S u = 0$. In other words $\Gamma S u$ has only terms that involve powers greater than 1 of $a_{i,j}$ for various $i,j$. Suppose that
\begin{align}
	\Gamma S u = a_{i,j}^2 f(x_1,x_2)
\end{align}
Then
\begin{align}
	u  = \frac{a_{i,j}^2}{S}\Gamma^{-1}(f)
\end{align}
So we here we can let
\begin{align}
		f = a_{k,l}^m x_1^rx_2^s t^n
\end{align}
and we find that
\begin{align}
	u =  \frac{a_{i,j}^2 a_{k,l}^m}{a_{r,s}S}x_1^r x_2^s t^n
\end{align}
Thus the kernel of $\Sigma$ is generated by these functions for all coefficients $i,j,k,l,r,s,n \in \ze$ and $m\ge 0$. This may perhaps prove useful later on. Now, if we look at eigenvalues $\lambda \neq 1$ we have to solve the equation
\begin{align}
	t \Sigma u = (1-\lambda) u
\end{align}
Because of the $t$, if $u$ is a finite order laurent series in $t$ with first non-zero coeficient $u_n$. We find that since $\Sigma$ does not affect the $t$ monomials
\begin{align}
	\Sigma u_n = \mu u_{n-1} = 0
\end{align}
and this forces all $u_k$ coefficients to be zero. Therefore $\Sigma$ has no other eigenvalues except 1 among the finite order power series in $t$.
%================================
\subsection{Generalizations}
%================================

The discussion lends itself nicely to generalization where now we have $d$ variables and operators $\Gamma_{i_1,\ldots,l_d}$, $R_{i_1,\ldots,i_d}$, etc defined analogously. The eigenpair discussion works out in the same way with eigenfunctions for $\lambda = 1$
\begin{align}
	u = \frac{a_{i_1,\ldots,i_d}^2 a_{j_1,\ldots,j_d}^m}{a_{k_1,\ldots,k_d}S}x_1^{r_1}\ldots x_d^{r_d} t^n
\end{align}
We can also constrain the SAWs on a part of $\ze^d$ with an extraction operator $E$ and mark special locations with an operator $L$. The functional equation is
\begin{align}
	(1- tLE\Sigma)W = a_{0,\ldots,0}
\end{align}
for arbitrary number of dimensions, with region constraints and marking of special locations.


\endinput
%================================================================
%================================================================
%================================================================
%================================================================
%================================================================

%================================
\section{Coefficient Extraction Equations}
%================================

The work so far is strongly hinting towards a theory of ordinary and partial coefficient extraction equations - OCEs and PCEs. Here we'll work out some of the theory. We start with a partial coefficient extraction operator of finite type
\begin{definition}$\mc{L}$ is a coefficient operator if
\begin{align}
	\mc{L} = \sum_{\substack{I\subseteq [d]\\ J \subseteq \ze^d}}a_{I,J}[x_I^J] = K + \mc{E}
\end{align}
where $K$ is the kernel of the operator, $\mc{E}$ is the pure extraction part
\begin{align}
	\mc{E}  = \sum_{\substack{I\subseteq [d],I\neq \emptyset\\ J \subseteq \ze^d}}a_{I,J}[x_I^J]
\end{align}
and all the coefficients are Laurent polynomials in the $x$ variables and perhaps some other variables. The operator is ordinary if $d = 1$ and partial if $d \ge 2$. The operator is either of monomial type where the coefficients are monomials or of polynomial type if one or more coefficients is a Laurent polynomial in $x_1,\ldots,x_d$.
\end{definition}
\begin{definition} An ordinary/partial coefficient equation (OCE/PCE) is one of the form $\mc{L} Q = q$ for some functions $Q$ and $q$ where $\mc{L}$ is an ordinary/partial coefficient operator. The equation is homogeneous if $q = 0$ and inhomogeneous if $q \neq 0$.
\end{definition}
The question is whether we can solve these kinds of equations in general and what their solutions may look like. In the previous section we found that if $\mc{L}$ is finite then the function $Q$ is D-finite in the $x$ variables while it may not necessarily be D-finite in the $t$ variable or any other variable (parameter) that appears in the coefficients.

\begin{example} The operator $\Gamma$ in this section is a partial coefficient operator with monomial coefficients.
\end{example}

\begin{example} We can transform the functional equation for a finite random walk model into a PCE. To do so we note that $K$ is the same in both forms. The part requiring some work is the one that has to do with the term $(1-\Gamma)SQ$ which is a pure extraction of $SQ$. We can expand $S$ into its monomials. Then $(1-\Gamma)S$ fully expands into monomials of the form
\begin{align}
	a_{I,J}\ext{x_I^J}x_A^B
\end{align}
we can pull terms inside $[x_I^J]$ from the right. In the end the $S$ is absorbed within the operator and istead of having monomial coefficients it now has polynomial coefficients. It may be useful to write an operator $\mc{L}$ in the form
\begin{align}
	\mc{L} = K + \Gamma'S
\end{align}
for some monomial operator $\Gamma'$.
\end{example}

%================================
\section{Connection with Fredholm Theory}
%================================
Suppose that we are given the equation in the form
\begin{align}
	KQ +t\Gamma'SQ = q
\end{align}
where
\begin{align}
	\Gamma' = \sum_{I,J}a_{I,J} x_I^J[x_I^J]
\end{align}
where we only have finite-sized $J$'s and $\emptyset \neq I\subseteq [d]$. From complex analysis we know that for a function $f(x)$ given by a power series and for $x$ within the region of convergence
\begin{align}
	[x^n]f(x) = f_n = \frac{1}{2\pi i}\int_\gamma \frac{f(z)}{z^{n+1}}dz
\end{align}
and
\begin{align}
	f(x) = \frac{1}{2\pi i}\int_\gamma \frac{f(z)}{z-x}dz
\end{align}
where in the first case $\gamma$ winds CCW around the origin within the region of convergence and in the second case $\gamma$ winds around $x$ within the region of convergence. The idea is to combine these two expressions for many dimensions (at least formally) to represent the operator $\mc{L}$ in integral form. We have for $|I| = n$
\begin{align}
	[x_I^J]f(x_1,\ldots,x_d)= \frac{1}{(2\pi i)^{n}}\int\cdots\int_\gamma \frac{f}{z_I^{J+1}}dz_I
\end{align}
where $f$ has $x_i = z_i$ for $i \in I$ and the rest are still $x$'s. If we use this representation we'll have different numbers of iterations of the integral for different coefficient. We might as well go all-out and write
\begin{align}
	[x_I^J]f(x_1,\ldots,x_d)= \frac{1}{(2\pi i)^d}\int\cdots\int_\gamma \frac{f}{z_I^{J+1}(z-x)_{I^c}}dz_{[d]}
\end{align}
where now $x_i = z_i$ in $f$ and 
\begin{align}
	(z-x)_P = \prod_{p\in P} (z_p-x_p)
\end{align}
This allows us to write The operator like this
\begin{align}
	(K+t\Gamma'S)Q&= \frac{1}{(2\pi i)^d}\int\cdots\int_\gamma \left(\frac{K}{(z-x)_{[d]}} +t S\sum_{I,J}\frac{a_{I,J}}{z_I^{J+1}(z-x)_{I^c}}\right)Q dz_{[d]}\\
	& = \frac{1}{(2\pi i)^d}\int\cdots\int_\gamma \frac{1}{(z-x)_{[d]}}\left(K +t S\sum_{I,J}\frac{a_{I,J}(z-x)_I}{z_I^{J+1}}\right)Q dz_{[d]}\\
	&= \frac{1}{(2\pi i)^d}\int\cdots\int_\gamma R_0Q dz_{[d]}
\end{align}
Now to turn this into a $d$-dimensional Fredholm problem we'll need to parametrize the integral. One way that might work is to set $z_k = r_k \exp(i\theta_k)$ and $x_k = s_k e^{i\phi_k}$. Since $dz_k = r_k i e^{i\theta_k}$ we have
\begin{align}
(K+t\Gamma'S)Q = \frac{r_{[d]}}{(2\pi)^d}\int_0^{2\pi}\cdots\int_0^{2\pi} e^{i\theta_1}\ldots e^{i\theta_d}R_0Q d\theta_{[d]} = q(\phi_1,\ldots,\phi_d)
\end{align}
so if we write
\begin{align}
R = \exp(i(\theta_1+\ldots + \theta_d))R_0(\phi_1,\ldots,\phi_d,\theta_1,\ldots,\theta_d)
\end{align}
we have the inhomogeneous Fredholm equation of the first kind in $d$ dimensions
\begin{align}
	\int_0^{2\pi}\cdots\int_0^{2\pi} RQ d\theta_1\ldots d\theta_d= g
\end{align}
where $R$ is the kernel and $g = (2\pi)^dq/r_{[d]}$ and $Q$ is in terms of $\theta_k$'s. Perhaps Fredholm theory can be applied in solving the random walk models in man dimensions. Perhaps this will also lead to a good inner product on the space where $Q$ lives.

%================================
\section{Reformulations of Problems}
%================================

\subsection{Using marking operators instead of extraction operators}
An initially promising idea is to use a marking operator instead of an extraction operator with the goal of removing those marked walks at a later stage. The added benefit is that we would be working with invertible operators throughout the majority of the analysis. The downside is that now the walks will not be constrained and $Q$ will no longer be a power series in the $x$ variables. We may get further issues. Instead of $E = 1- \ext{\ol{x}}$ we take $\tilde E = 1 + (a_{-1}-1)\ext{x}$ with the idea that if we set $a_{-1} = 0$ we'll obtain $E$. To this end the functional equation is
\begin{align}
	(1-tL\tilde E S)\tilde Q = 1
\end{align}
which after inversion turns into
\begin{align}
	(\tilde E^{-1} L^{-1} -tS) \tilde Q = \tilde E^{-1} L^{-1} 1 = \alpha
\end{align}
One benefit of this is that we don't have a mixing between the $E$, $L$ operators and $S$. The downside is we have more $Q$ extractions on their own. We can write this in kernel form as
\begin{align}
	K \tilde Q = \alpha + (1 - L^{-1}\tilde E^{-1}) \tilde Q
\end{align}
where the operator $1- L^{-1}\tilde E^{-1}$ is composed entirely of coefficient extractions.
\subsection{Considering a symmetrical super-problem}
An extension of the previous idea is to symmetrize the problem by also marking parts on the other side of the origin. Instead of just marking walks on the axes and on lines where one coordinate is $-1$ we can also mark walks on lines with one coordinate $+1$. This leads to a symmetrical arrangement of marking operations on both sides of the origin. We have the equation
\begin{align}
	(1- tLS)\tilde Q = (1- t L_{-1}L_0L_1 S)\tilde Q = 1
\end{align}
We can skip the marking of the axes with $L_0$ if we don't wish to investigate wall interactions. In general we may mark more locations in space so we can combine all of these operators into one called $L$. Then we have
\begin{align}
	(1- tLS)\tilde Q = 1
\end{align}
or in kernel forms
\begin{align}
	K\tilde Q = L^{-1}(1) + (1- L^{-1})\tilde Q
\end{align}

\endinput
\begin{example} Here is an example that may hint at the flavour of the OCE/PCE problems. Suppose we start with the fibonacci recurrence relation
\begin{align}
F_{n+2} =F_{n+1}+F_n
\end{align}
with some initial conditions for $f_0$ and $f_1$. We can perform the standard generating function solution approach via summing all equations. Let
\begin{align}
	F(x) = \sum_{n=0}^\infty F_n x^n
\end{align}
be the generating function for the fibonacci-like numbers. Note that $[x^n] F(x) = F_n$. We have
\begin{align}
	F_{n+2} &= F_{n+1}+F_n\\
	F_{n+2} x^{n+2}&= F_{n+1}x^{n+2}+F_nx^{n+2}\\
	\sum_{n=0}^{\infty} F_{n+2} x^{n+2}&= x\sum_{n=0}^{\infty}F_{n+1}x^{n+1}+x^2\sum_{n=0}^{\infty}F_nx^n\\
	F - F_1 x - F_0 &= x (F-F_0) + x^2 F
\end{align}
Pulling like terms together and putting everything on one side gives
\begin{align}
	(1-x - x^2) F - x F_1 + (x-1) F_0 = 0
\end{align}
This is a homogeneous OCE with ordinary coefficient operator$\mc{L}$ (Fibonacci operator?) with kernel $K = 1 - x - x^2$ and polynomial coefficients
\begin{align}
	\mc{L} = (1-x-x^2) + (x-1)[x^0]-x[x] = K + \mc{E}
\end{align}
The typical solution to the problem specifies initial conditions $F_0$ and $F_1$ and solves the problem as
\begin{align}
	F(x) = \frac{(x-1)F_0-x F_1}{1-x-x^2} = \frac{{\mc{E}}(F)}{K}
\end{align}
That is, it relies on the initial condition for $\mc{E}$ in order to solve the equation. We also see that if we don't specify the initial conditions we don't get a unique solution. In the random walks models we have unique solutions because of the iterative nature of the problems. With these points in mind we may try the vanishing kernel method by supplying a root of $K$ namely $x = \sigma$. These roots are
\begin{align}
\sigma_1 = -\varphi, \sigma_2 = \ol{\varphi}
\end{align}
where $\varphi$ is the golden ratio. Then plugging either one of these into the OCE leads to
\begin{align}
(\sigma-1)F_0 = \sigma F_1
\end{align}
\begin{align}
	 F_1 &= (1-\ol{\sigma}_1)F_0 = \varphi F_0\\
	 \text{ or } F_1 &=(1-\ol{\sigma}_2)F_0 = -\ol{\varphi} F_0
\end{align}
That is, if we want the vanishing kernel method to work we must have that $F_1$ is related to $F_0$ directly via the golden ratio. By induction we can show that $F_n = \varphi^n F_0$ in the first case and $F_n = (-\ol{\varphi})^n F_0$ in the second case.

What about the persistent kernel method? We want to apply a transformation $y = y(x)$ to $K$ that leaves it the same. That is
\begin{align}
	K(y) = 1 - y - y^2 = K(x) = 1 - x - x^2
\end{align}
Collecting terms on one side yields
\begin{align}
	(y-x)(1+x+y) = 0
\end{align}
That is, the group $G$ that preserves the kernel is composed of the identity transformation and the the transformation $g(x) = -x-1$. We want a $F(-1-x)$ to converge which may not necessarily be true. Instead we multiply the whole functional equation by $\ol{x}$ to find
\begin{align}
\ol{x} \mc{L} = (\ol{x}-1-x) + (1-\ol{x})[x^0]-[x^1]
\end{align}
This new kernel $K$ is nicer. We have
\begin{align}
	K = \ol{x}-1-x = \ol{y}-1-y
\end{align}
rearranging this equality gives
\begin{align}
	(y-x)(1+\ol{x}\ol{y}) = 0
\end{align}
and thus the group $G$ is generated by the identity transformation and $g(x) = -\ol{x}$ which is of order 2, so $|G| = 2$. We have the two equation
\begin{align}
KF +(1-\ol{x})F_0 - F_1 &= 0\\
KgF + (1+x)F_0  - F_1 &=0
\end{align}
Furthermore emulating the orbit method we have
\begin{align}
	(1-g)xKF = K(xF+\ol{x}gF)  = (x+\ol{x})(F_1-F_0)
\end{align}
If $F_1 \neq F_0$ we can write
\begin{align}
	xF+\ol{x}gF &= \frac{(x+\ol{x})}{K}(F_1-F_0) \\
	&= \frac{1+x^2}{1-x-x^2}(F_1-F_0) \\
	&= (1+x+3x^2+O(x^3))(F_1-F_0)
\end{align}
Extracting the $[x]$ coefficient yields
\begin{align}
	F_0 +[x^2]F(-\ol{x}) = F_0 = F_1 - F_0
\end{align}
so $F_1= 2 F_0$ and we find
\begin{align}
	xF+\ol{x}F(-\ol{x}) = F_0 \frac{1+x^2}{1-x-x^2}
\end{align}
Taking the positive part of this yields
\begin{align}
	xF= F_0 \frac{1+x^2}{1-x-x^2}
\end{align}
\end{example}
I think there is something wrong with the example but that's just me. We can demonstrate the D-finiteness by taking the second derivative of $xKF$
\begin{align}
	\del_x^2((1-x - x^2) F - x F_1 + (x-1) F_0) = \del_x^2(xKF) = 0
\end{align}